\section{Plan de validación}

El plan de validación del sistema portátil de limpieza CIP se estructura a partir de la matriz de
requisitos definida en la Sección~3.1 (Tabla~6) y de las decisiones de diseño de control y
comunicaciones presentadas en las Secciones~5 y~6. Para el segundo corte se prevé ejecutar como mínimo tres
ciclos completos del proceso CIP (preenjuague, lavado y enjuague final) y una prueba continua,
 evaluando al menos tres valores de consigna de temperatura dentro del rango seguro
(35\,°C, 40\,°C y 50\,°C) y aplicando una perturbación controlada mediante la adición de fluido frío
durante la fase de lavado. En este corte se presenta únicamente el plan; los resultados experimentales
se reportarán en el segundo corte, una vez ejecutadas las pruebas sobre el prototipo integrado.

Cada prueba se identifica con un código, que se relaciona con uno o más requisitos de la
Tabla~6 y define el procedimiento, el instrumento de referencia y el criterio de aceptación que permitirá
declarar el requisito como cumplido.

\begin{longtable}{p{1.3cm} p{1.4cm} p{4.3cm} p{2.6cm} p{3.4cm} p{2.4cm}}
\caption{Matriz requisito--prueba--evidencia.}
\label{tab:matriz-validacion} \\
\toprule
\textbf{ID} & \textbf{Requisito} & \textbf{Procedimiento} & \textbf{Instrumento} & \textbf{Criterio} & \textbf{Evidencia} \\
\midrule
\endfirsthead

\multicolumn{6}{c}{\tablename\ \thetable{} -- continuación} \\
\toprule
\textbf{ID} & \textbf{Requisito} & \textbf{Procedimiento} & \textbf{Instrumento} & \textbf{Criterio} & \textbf{Evidencia} \\
\midrule
\endhead

\midrule
\multicolumn{6}{r}{Continúa en la siguiente página} \\
\endfoot

\bottomrule
\endlastfoot

VAL-01 & REQ01 & Ejecutar el ciclo completo automático (preenjuague, lavado, enjuague) desde el HMI, sin intervención manual, repetido en 3 pruebas consecutivas & Cronómetro, registro del PLC & 3/3 ejecuciones completas en la secuencia correcta & Trazas de ciclo exportadas del HMI \\

VAL-02 & REQ02 & Verificar la apertura y cierre de SV-101, SV-102 y SV-103 en cada fase, confirmando ausencia de traslapes entre etapas incompatibles & Registro de estados digitales del PLC & Secuencia correcta sin traslapes en 3 pruebas consecutivas & Trazas de variables digitales \\

VAL-03 & REQ13 & Consignar SP = 50\,°C (prelavado) y SP = 40\,°C (lavado) y registrar la evolución de la temperatura hasta alcanzar el estado estable & TT-101 (PT100 0 SM1231), termómetro patrón & Error en estado estable $ \leq 2\,°C $ respecto al SP & Gráfica $T(t)$ vs. SP y datos exportados \\

VAL-04 & REQ14 & Forzar, de forma controlada, una condición de sobre temperatura y verificar el corte del calentador E-101 (vía SSR)  & Registro de eventos del PLC, cronómetro & Corte del SSR y aviso activado al superar 60\,°C & Registro de evento y captura de HMI \\

VAL-05 & REQ09 & Dosificar, mediante la bomba peristáltica P-101, 3 repeticiones de un mismo tiempo de activación y medir el volumen de jabón entregado en cada una & Probeta graduada, cronómetro & Variación $\leq 10\,\%$ entre las 3 mediciones & Registro fotográfico y tabla de volúmenes \\

VAL-06 & REQ06 & Ejecutar un ciclo completo e inspeccionar los tanques TK-101, TK-102 y TK-103 y sus conexiones (SV-101, SV-102, SV-103, P-101, P-102) en busca de fugas & Inspección visual directa & Cero fugas visibles durante la prueba & Registro fotográfico \\

VAL-07 & REQ06 & Medir el volumen de agua suministrado al ciclo (línea L-2) y el volumen recuperado en el tanque de almacenamiento TK-103 (línea L-4) al finalizar & Recipiente graduado o báscula & Porcentaje de recuperación calculado y reportado & Registro de volúmenes y cálculo de \% \\

VAL-08 & REQ18 & Accionar el paro de emergencia durante la ejecución del ciclo y medir el tiempo de corte de bomba, calentador y dosificador & Cronómetro, registro de eventos del PLC & Corte de los actuadores principales en $\leq 1\,\text{s}$ & Captura de señal o video \\

VAL-09 & REQ19 & Simular las 3 condiciones anómalas previstas (nivel bajo, sobretemperatura, falla de bomba) y verificar la activación de la alarma visual y sonora & Inspección directa & Aviso activado correctamente en las 3 condiciones simuladas & Registro fotográfico o video \\

VAL-10 & REQ15, REQ16 & Verificar, durante una prueba, la actualización de temperatura, nivel y estado del ciclo en el HMI y en la plataforma de supervisión IoT & Captura de tráfico de red, trazas de la plataforma & $\geq 99\,\%$ de los registros esperados recibidos, con periodo $\leq 30\,\text{s}$ & Trazas de red y capturas de la plataforma \\

VAL-11 & REQ11 & Comparar la lectura del sensor de nivel LIT-101, instalado en el tanque de mezcla TK-102, contra una medición manual de referencia en al menos 3 niveles distintos & LIT-101, regla o marca graduada & Error $\leq 5\,\%$ del rango medido en las 3 mediciones & Tabla de lecturas y registro fotográfico \\

VAL-12 & REQ01 (PMV, Sección~4.5) & Activar el motorreductor M-101 y el agitador AG-101 durante la fase de preparación de la mezcla en TK-102, y verificar la homogeneidad del fluido tras un tiempo de agitación definido, repitiendo el ensayo 3 veces & Inspección visual directa, cronómetro & Mezcla visualmente homogénea (sin zonas de concentración diferenciada) en las 3 pruebas, dentro del tiempo de agitación definido & Registro fotográfico o video \\

\end{longtable}

Adicionalmente, para cada consigna evaluada se comparará la respuesta experimental de temperatura con
la respuesta simulada del modelo obtenido en la Sección~5.5, con el fin de analizar las diferencias,
las saturaciones del actuador y las limitaciones del modelo de primer orden adoptado.

\subsection{Métricas previstas}

Las métricas que se calcularán a partir de los datos registrados durante las pruebas del segundo corte se
resumen en la Tabla~20. El método de cálculo de cada métrica se apoya en los datos exportados desde el
PLC, el HMI y la plataforma de supervisión IoT, evitando el uso de valores simulados como si fueran
mediciones reales.

\begin{longtable}{p{2.6cm} p{4.6cm} p{2.0cm} p{5.0cm}}
\caption{Métricas previstas para validación.}
\label{tab:metricas-validacion} \\
\toprule
\textbf{Área} & \textbf{Métrica} & \textbf{Unidad} & \textbf{Método de cálculo} \\
\midrule
\endfirsthead

\multicolumn{4}{c}{\tablename\ \thetable{} -- continuación} \\
\toprule
\textbf{Área} & \textbf{Métrica} & \textbf{Unidad} & \textbf{Método de cálculo} \\
\midrule
\endhead

\midrule
\multicolumn{4}{r}{Continúa en la siguiente página} \\
\endfoot

\bottomrule
\endlastfoot

Control & Error estacionario, sobreimpulso y tiempo de establecimiento & \%, °C &
\begin{minipage}[t]{5.0cm}
\vspace{2pt}
$\displaystyle e_{ss} = \dfrac{|SP - T_{estable}|}{SP}\times 100$

\vspace{6pt}
$\displaystyle M_p = \dfrac{T_{max}-SP}{SP}\times 100$

\vspace{6pt}
Calculados a partir de la curva $T(t)$ exportada del PLC/HMI.
\end{minipage} \\[4pt]
Comunicaciones & Latencia, pérdida de mensajes y tiempo de reconexión & ms, \%, s & Captura de tráfico de red con marcas de tiempo; conteo de paquetes esperados frente a recibidos \\
Dosificación & Error absoluto y relativo del volumen dosificado & mL, \% & Comparación entre el volumen programado (por tiempo de activación) y el volumen medido en probeta, en 3 repeticiones \\
Recuperación & Porcentaje de recuperación de agua & \% & Relación entre el volumen recuperado en TK-103 y el volumen total suministrado al ciclo \\
Instrumentación & Error de medición de nivel (LIT-101) & \% & Comparación entre la lectura del sensor y la medición manual de referencia en TK-102, en al menos 3 niveles \\
Agitación & Tiempo de homogeneización de la mezcla (M-101/AG-101) & min & Cronometraje desde el inicio de la agitación hasta lograr una mezcla visualmente homogénea, en 3 repeticiones \\
Seguridad & Tiempo de respuesta del paro de emergencia y de las alarmas & s & Cronometraje o captura de eventos del PLC desde la activación hasta el corte o la señal de alarma \\

\end{longtable}

Estas métricas permitirán, en el segundo corte, contrastar el desempeño experimental del sistema frente a las
especificaciones preliminares definidas en la Tabla~10 y frente a los criterios mínimos de aceptación
establecidos para el proyecto.